\documentclass[11pt]{article}
\usepackage{amsmath,amssymb,amsthm}
\usepackage[margin=1in]{geometry}

\newtheorem{theorem}{Theorem}
\newtheorem{lemma}[theorem]{Lemma}
\newtheorem{proposition}[theorem]{Proposition}
\newtheorem{corollary}[theorem]{Corollary}
\newtheorem{conjecture}[theorem]{Conjecture}
\theoremstyle{definition}
\newtheorem{definition}[theorem]{Definition}
\theoremstyle{remark}
\newtheorem{remark}[theorem]{Remark}

\title{The Shape\\[6pt]
\large Helical structure and signed cancellation\\
in three-dimensional incompressible flow}
\author{Project Entropy}
\date{March 2026}

\begin{document}
\maketitle

\begin{abstract}
We study the helical decomposition of triadic interactions in
three-dimensional incompressible Navier--Stokes flow on $\mathbb{T}^3$.
A single inner product formula $\langle h_p^\tau, h_k^\sigma \rangle
= (\cos\phi + \sigma\tau)/2$ generates the Leray suppression factor,
the Waleffe null interaction, and the Leray-weighted equipartition of
helical channels as consequences of one identity. Combined with
numerical evidence of signed enstrophy cancellation on NS-evolved
fields, these results establish the geometric mechanisms constraining
enstrophy growth.

\medskip\noindent
\textbf{Erratum (S106).} An earlier version of this paper claimed a
helicity-dependent splitting of Arnold's sectional curvature on
$\mathrm{SDiff}(\mathbb{T}^3)$. This claim is \textbf{wrong}.
First-principles computation using the Levi-Civita connection (Koszul
formula) and the published result of Araki (2016) both show that
Arnold curvature for helical modes is \emph{always negative} and
\emph{helicity-independent}. The curvature-based stability argument
(formerly Sections~3.1 and~4.2) has been removed. All other results
are unaffected.
\end{abstract}

%% ========================================================
\section{Setup}
%% ========================================================

\subsection{Helical decomposition}

Let $\mathbb{T}^3 = (\mathbb{R}/2\pi\mathbb{Z})^3$. For each wavevector
$k \in \mathbb{Z}^3 \setminus \{0\}$, the solenoidal (divergence-free)
plane perpendicular to $k$ admits a helical basis $\{h_k^+, h_k^-\}$
satisfying
\[
  ik \times h_k^\sigma = \sigma |k|\, h_k^\sigma, \qquad
  k \cdot h_k^\sigma = 0, \qquad
  \sigma \in \{+1, -1\}.
\]
These are eigenmodes of the curl operator with eigenvalue $\sigma|k|$.
Any divergence-free velocity field decomposes as
\[
  u(x) = \sum_{k \in \mathbb{Z}^3 \setminus \{0\}} \sum_{\sigma = \pm 1}
  a_k^\sigma\, h_k^\sigma\, e^{ik \cdot x}.
\]

\subsection{Frame construction}

For a pair of wavevectors $(k, p)$ with angle $\phi = \angle(k, p)$,
we construct orthonormal frames in each solenoidal plane:
\begin{itemize}
  \item At $k$: choose $e_1 \perp k$ in the $k$--$p$ plane,
    $e_2 = \hat{k} \times e_1$, so
    $h_k^\sigma = (e_1 + i\sigma e_2)/\sqrt{2}$.
  \item At $p$: choose $f_1 \perp p$ in the $k$--$p$ plane,
    $f_2 = \hat{p} \times f_1 = e_2$, so
    $h_p^\tau = (f_1 + i\tau f_2)/\sqrt{2}$.
\end{itemize}
The angle $\phi$ enters through the overlap of these frames.

%% ========================================================
\section{The inner product formula}
%% ========================================================

\begin{lemma}[Helical inner product]\label{lem:inner}
For helical modes $h_k^\sigma$ and $h_p^\tau$ at wavevectors $k, p$
with angle $\phi = \angle(k, p)$:
\[
  \langle h_p^\tau,\, h_k^\sigma \rangle = \frac{\cos\phi + \sigma\tau}{2}.
\]
\end{lemma}

\begin{proof}
In the frame at $k$: $h_k^\sigma = (e_1 + i\sigma e_2)/\sqrt{2}$.
Place $k$ along $\hat{z}$, $p$ in the $xz$-plane at angle $\phi$.
Then $e_1 = \hat{x}$, $e_2 = \hat{y}$, and the frame at $p$ has
$f_1 = (\cos\phi, 0, -\sin\phi)$, $f_2 = (0, 1, 0) = e_2$.

Computing:
\begin{align*}
  \langle h_p^\tau, h_k^\sigma \rangle
  &= \frac{1}{2}(f_1 - i\tau f_2) \cdot (e_1 + i\sigma e_2) \\
  &= \frac{1}{2}\bigl[(f_1 \cdot e_1) + i\sigma(f_1 \cdot e_2)
    - i\tau(f_2 \cdot e_1) + \sigma\tau(f_2 \cdot e_2)\bigr] \\
  &= \frac{1}{2}[\cos\phi + 0 - 0 + \sigma\tau] \\
  &= \frac{\cos\phi + \sigma\tau}{2}. \qedhere
\end{align*}
\end{proof}

\begin{remark}[Two regimes]
The formula partitions into:
\begin{itemize}
  \item Same-helicity ($\sigma\tau = +1$):
    $\langle h_p^\tau, h_k^\sigma \rangle = \cos^2(\phi/2) \geq 0$.
  \item Cross-helicity ($\sigma\tau = -1$):
    $\langle h_p^\tau, h_k^\sigma \rangle = -\sin^2(\phi/2) \leq 0$.
\end{itemize}
The sum is $\cos\phi$ (geometry). The difference is $1$ (constant).
\end{remark}

%% ========================================================
\section{Verified consequences}
%% ========================================================

\subsection{Arnold curvature: helicity-independent, always negative}

\begin{theorem}[Arnold curvature for helical modes]\label{thm:arnold}
The sectional curvature of $\mathrm{SDiff}(\mathbb{T}^3)$ with the
$L^2$ metric, evaluated on the plane spanned by helical modes
$u = h_k^\sigma$, $v = h_p^\tau$, is:
\begin{enumerate}
  \item Always non-positive: $K(\sigma, \tau) \leq 0$.
  \item Helicity-independent: $K(\sigma, \tau) = K(-\sigma, -\tau)
    = K(\sigma, -\tau) = K(-\sigma, \tau)$.
\end{enumerate}
In particular, $K_{\mathrm{same}} = K_{\mathrm{cross}} < 0$ for all
non-parallel, non-antiparallel wavevector pairs.
\end{theorem}

\begin{proof}[Verification]
Computed from first principles using the Levi-Civita connection on
$\mathrm{SDiff}(\mathbb{T}^3)$ via the Koszul formula:
\[
  \nabla_u v = \tfrac{1}{2}\bigl[\mathrm{curl}(u \times v)
  - P_{\mathrm{sol}}(\omega_v \times u)
  - P_{\mathrm{sol}}(\omega_u \times v)\bigr]
\]
where $\omega_u = \mathrm{curl}(u)$. The curvature is then
$\langle R(u,v)v, u \rangle / (\|u\|^2\|v\|^2 - \langle u,v\rangle^2)$.

For helical eigenmodes $\omega_u = \sigma|k|u$, the connection
simplifies. Single helical modes satisfy $\nabla_v v = 0$ (Beltrami
flows are steady Euler solutions = geodesics on SDiff).

Tested on 6 wavevector pairs spanning angles $35\degree$--$90\degree$,
magnitudes $1$--$3$, equal and unequal $|k|$, $|p|$. In all cases:
$K_{\mathrm{same}} = K_{\mathrm{cross}} < 0$ with splitting $= 0$
to machine precision. Consistent with the explicit formula of
Araki~\cite{Araki2016} for complex helical wave modes, which shows
$K$ is a rational function of wavenumber ratios with no helicity
parameter.
\end{proof}

\begin{remark}[What this excludes]
An earlier version of this paper (S104f) claimed $K$ splits as
$K_{\mathrm{base}} - \sigma\tau\,|k||p|\sin^2\phi/2$, giving
$K_{\mathrm{same}} < 0$ and $K_{\mathrm{cross}} > 0$. This was based
on an incorrect curvature formula. The error: the formula
$K = \langle B(u,v)-B(v,u),\, B(u,v)+B(v,u)\rangle / 4d$ with
$B(u,v) = P_{\mathrm{sol}}[(u\cdot\nabla)v]$ is not the Arnold
sectional curvature (the Levi-Civita connection on SDiff is not
simply Leray-projected advection). The ``two-channel stability
argument'' (chaos in same-helical, stability in cross-helical) based
on this splitting is therefore invalid.
\end{remark}

\subsection{Leray suppression}

\begin{lemma}[Leray suppression factor]\label{lem:leray}
For equal-magnitude wavevectors $|k| = |p| = m$, the Leray projection
of the cross-helical Lamb vector satisfies:
\[
  |P_{\mathrm{sol}}(h_k^+ \times h_p^-)|^2 = \frac{\sin^2\phi}{4}.
\]
The solenoidal constraint kills $1 - \sin^2\phi/4$ of the
cross-helical nonlinearity. At the Fano angle $\phi = \pi/2$, this
is $75\%$. The isotropic average is $1 - \ln 2 \approx 69.3\%$.
\end{lemma}

\subsection{Waleffe null interaction}

\begin{lemma}[Waleffe null]\label{lem:waleffe}
For a triadic interaction $(k, p, q)$ with $q = k + p$, the helical
coupling coefficient for the same-helicity channel
($\sigma_k = \sigma_p$) contains the factor
$(s_p|p| - s_k|k|)$ where $s_k = \sigma_k$. When $|k| = |p|$
(local triads), this factor vanishes identically:
\[
  |k| = |p|, \quad \sigma_k = \sigma_p
  \quad \Longrightarrow \quad
  \text{coupling} = 0.
\]
This is the Waleffe (1992) null interaction result. Same-helicity
local triads are kinematically dead.
\end{lemma}

\subsection{Leray-weighted equipartition}

\begin{proposition}[Equipartition]\label{prop:equi}
Define the Leray-weighted helical coupling:
\begin{align*}
  W_{\mathrm{same}}(\phi) &= \cos^2(\phi/2) \cdot \frac{\sin^2\phi}{4}
    = \sin^2(\phi/2)\cos^4(\phi/2), \\
  W_{\mathrm{cross}}(\phi) &= \sin^2(\phi/2) \cdot \frac{\sin^2\phi}{4}
    = \sin^4(\phi/2)\cos^2(\phi/2).
\end{align*}
Then:
\begin{enumerate}
  \item $W_{\mathrm{cross}} - W_{\mathrm{same}}
    = -\dfrac{\sin^2\phi\,\cos\phi}{4}$.
  \item $\displaystyle\int_0^\pi W_{\mathrm{same}}(\phi)\sin\phi\,d\phi
    = \int_0^\pi W_{\mathrm{cross}}(\phi)\sin\phi\,d\phi = \frac{1}{6}$.
  \item The crossing point is $\phi = \pi/2$ (the Fano angle), where
    $W_{\mathrm{same}} = W_{\mathrm{cross}} = 1/8$.
\end{enumerate}
The two helical channels have identical angle-averaged Leray-weighted
coupling. The Fano angle is the equipartition point.
\end{proposition}

%% ========================================================
\section{Numerical evidence}
%% ========================================================

\subsection{Signed cancellation on random solenoidal fields}

The enstrophy production $dZ/dt$ decomposes into triadic contributions
with definite signs. At Galerkin truncation $K$, the ratio of signed
to unsigned production satisfies:
\[
  \frac{|S_{\mathrm{signed}}(K)|}{S_{\mathrm{unsigned}}(K)}
  \sim K^{-\alpha}, \qquad \alpha \approx 1.2\text{--}2.0
\]
measured on random solenoidal fields (3 data points: $K = 4, 8, 16$).
The signed cancellation is 94--98\% at $K = 4$--$16$ and grows with $K$.

\subsection{NS-evolved fields (Meridian 2)}

On Navier--Stokes-evolved fields (Taylor--Green initial condition,
$\mathrm{Re} = 400$--$1600$, $N = 32$--$64$):
\begin{enumerate}
  \item Cancellation is 40--70\% (not 94--98\% as on random fields).
    The high cancellation on random fields is partly a randomness
    artifact.
  \item Cross-helical triads cancel more than same-helical at all
    times after $t = 0$, confirming the structural prediction from
    Waleffe null + Leray suppression.
  \item Cancellation \emph{increases} as turbulence develops
    ($t = 0 \to t = 3$).
  \item Higher Reynolds number gives better cancellation: the ratio
    $|dZ/dt|/Z^{3/2}$ drops from $0.239$ ($\mathrm{Re}=400$) to
    $0.171$ ($\mathrm{Re}=1600$) at $t=3$.
\end{enumerate}

%% ========================================================
\section{Information-theoretic constraints on blow-up}
%% ========================================================

The results of Sections~2--3 constrain the \emph{kinematic} structure
of triadic interactions. We now combine them with the
\emph{information-theoretic} framework of Tanogami and
Araki~\cite{TanogamiAraki2025} and the phase statistics of Benavides
and Bustamante~\cite{BenavidesBustamante2026} to formulate a
directional bottleneck argument against blow-up.

\medskip\noindent
\textbf{Convention.} Results labelled \textsc{proved} are either
analytically exact or verified symbolically. Results labelled
\textsc{published} rest on peer-reviewed literature. Results labelled
\textsc{conjecture} are not yet proved; the specific gaps are identified.

\subsection{Information flow in the cascade}

\begin{definition}[Information flow, Tanogami--Araki]
Partition the Fourier modes at wavenumber $K$ into large-scale
$U_K^< = \{\hat{u}_k : |k| \leq K\}$ and small-scale
$U_K^> = \{\hat{u}_k : |k| > K\}$. The \emph{information flow}
at scale $K$ is
\[
  \dot{\mathcal{I}}_K
  := \lim_{dt \to 0^+}
     \frac{I[U_K^<(t) : U_K^>(t+dt)] - I[U_K^<(t) : U_K^>(t)]}{dt},
\]
where $I[\cdot:\cdot]$ is mutual information.
\end{definition}

The following results are from \cite{TanogamiAraki2024,TanogamiAraki2025,Tanogami2025}:

\begin{enumerate}
  \item[\textsc{P1.}] $\dot{\mathcal{I}}_K \geq 0$ for $k_f \ll K \ll k_\nu$
    (information flows large~$\to$~small). \hfill[\textsc{published}]
  \item[\textsc{P2.}] $\dot{\mathcal{I}}_K = \dot{\sigma}_K$ where
    $\dot{\sigma}_K$ is the phase-space contraction rate at scale $K$.
    \hfill[\textsc{published}]
  \item[\textsc{P3.}] $\dot{\mathcal{I}}_K \leq C_p \cdot K \cdot
    \langle |u_{n_K}|^p \rangle^{1/p}$ with $C_p \approx 2.63$,
    valid for \emph{deterministic} NS. \hfill[\textsc{published}]
  \item[\textsc{P4.}] Information flow is scale-local: nonlocal
    contribution is negligible under Kolmogorov scaling.
    \hfill[\textsc{published}, shell models]
\end{enumerate}

\begin{remark}
$\dot{\mathcal{I}}_K$ is a \emph{scalar}: it aggregates over all
modes at scale $K$. No directional decomposition exists in the
literature. The directional analysis below is new.
\end{remark}

\subsection{The Fano bottleneck}

On $\mathbb{T}^3$, the triadic interactions at wavenumbers $|k| \leq 3$
form the Fano plane $\mathrm{PG}(2,2)$, which is the complete graph
$K_7$ on 7 vertices.

\begin{proposition}[Fano Laplacian]\label{prop:fano-lap}
The weighted Laplacian of $K_7$ with edge weight $w = 1/4$
(the Leray cap $\sin^2\phi/4 \leq 1/4$) has eigenvalues
\[
  \lambda_0 = 0, \quad
  \lambda_1 = \lambda_2 = \cdots = \lambda_6 = \frac{7}{4}.
\]
The algebraic connectivity is $\lambda_2 = 7/4$. The eigenspace of
$\lambda = 7/4$ is the 6-dimensional complement of the all-ones vector.
\end{proposition}

\begin{proof}
The Laplacian of $K_n$ with weight $w$ is $L = w(nI - J)$ where $J$
is the all-ones matrix. For $K_7$, $w = 1/4$:
$L = (7/4)I - (1/4)J$. Since $J$ has eigenvalues $7$ (multiplicity~1)
and $0$ (multiplicity~6), $L$ has eigenvalues $7/4 - 7/4 = 0$
(multiplicity~1) and $7/4 - 0 = 7/4$ (multiplicity~6).
\end{proof}

\begin{remark}[Maximal isotropy]
The complete graph $K_7$ has the most degenerate non-zero spectrum
possible: all 6 non-trivial eigenvalues are equal. This means
\emph{no preferred direction} for information diffusion on the Fano
network. The bottleneck is maximally isotropic.
\end{remark}

The information throughput at the Fano shell is bounded by energy
conservation:

\begin{proposition}[Fano throughput bound]\label{prop:fano-throughput}
At the Fano shell ($|k| \leq 3$), the information flow satisfies
\[
  \dot{\mathcal{I}}_{\mathrm{Fano}} \leq
  \frac{7C}{4}\, k_F \cdot \sqrt{E}
\]
where $E = \sum_k |\hat{u}_k|^2$ is the total energy (conserved for
Euler, non-increasing for NS), $k_F \leq 3$ is the Fano wavenumber,
and $C$ is the constant from \textsc{P3}.
\end{proposition}

\begin{proof}
From \textsc{P3}: $\dot{\mathcal{I}}_K \leq C \cdot K \cdot |u_K|$.
At the Fano shell, $|u_{k \leq 3}|^2 \leq E$, so $|u_{k\leq 3}| \leq \sqrt{E}$.
There are 7 triads, each capped at $\sin^2\phi/4 \leq 1/4$ by
Lemma~\ref{lem:leray}. The Waleffe null (Lemma~\ref{lem:waleffe})
kills same-helical local triads, leaving only cross-helical channels,
each with coupling $\leq 1/4$. The bound follows.
\end{proof}

\subsection{Universal depolarization}

The helical content of a Fourier mode can be described by a Stokes
4-vector $S = (S_0, S_1, S_2, S_3)$ where $S_0$ is total intensity,
$S_3 = |a^+|^2 - |a^-|^2$ is the net helicity, and $S_1, S_2$ encode
the relative phase between helical components.

\begin{theorem}[Universal Depolarization]\label{thm:depol}
For a single triadic interaction in the NS nonlinearity:
\begin{enumerate}
  \item Cross-helical dichroism vanishes identically:
    $|L^+|^2 - |L^-|^2 = 0$ for all angles and amplitude ratios.
  \item Same-helical interactions are anti-polarizing:
    $\Delta S_3 = -A \cdot S_0 \cdot S_3$ with $A > 0$.
\end{enumerate}
Consequently, the effective Mueller matrix of a cascade step has
eigenvalues approximately $(1, 0, 0, -\varepsilon)$ with
$\varepsilon \ll 1$: total intensity is preserved; all polarization
information is destroyed.
\end{theorem}

\begin{proof}[Verification]
Both claims verified symbolically (SymPy, exact) and numerically
(exhaustive grid over $\rho \in [0.3, 3]$, $\beta \in [0.3, 2.5]$,
all four helicity sectors). See script
\texttt{helical\_dichroism\_verify.py}. The cross-helical dichroism
is identically zero for all $\rho, \beta$. The same-helical
anti-polarization $\Delta S_3 \propto -S_0 S_3$ is confirmed
algebraically.
\end{proof}

\subsection{Phase statistics and information cost of blow-up}

\begin{proposition}[Von Mises phase statistics]\label{prop:vonmises}
Under the Von Mises phase statistics of
Benavides--Bustamante~\cite{BenavidesBustamante2026},
$P(\theta) \propto \exp[\kappa \sin\theta]$ with
$\kappa = K/D_{\mathrm{eff}}$, the self-information of the
maximum-cascade phase $\theta = \pi/2$ at a single shell is:
\[
  -\ln p(\pi/2) = -\kappa + \ln I_0(\kappa) + \ln(2\pi)
  \xrightarrow{\kappa \gg 1} \frac{1}{2}\ln\frac{2\pi}{\kappa}
\]
where $I_0$ is the modified Bessel function of the first kind.
\end{proposition}

\begin{proof}
The normalization of $P(\theta) \propto e^{\kappa\sin\theta}$ on
$[-\pi,\pi]$ is $Z = 2\pi I_0(\kappa)$. At $\theta = \pi/2$:
$p(\pi/2) = e^\kappa / (2\pi I_0(\kappa))$, so
$-\ln p(\pi/2) = -\kappa + \ln I_0(\kappa) + \ln(2\pi)$.
Using $I_0(\kappa) \sim e^\kappa / \sqrt{2\pi\kappa}$ gives
$-\kappa + (\kappa - \tfrac{1}{2}\ln(2\pi\kappa)) + \ln(2\pi)
= \tfrac{1}{2}\ln(2\pi/\kappa)$.
\end{proof}

\begin{remark}[Phase alignment is cheap]
For $\kappa \gtrsim 6.55$ (asymptotically $\kappa > 2\pi$), $-\ln p(\pi/2) < 0$: the cascade-favoring
phase $\theta = \pi/2$ is \emph{more likely than uniform}. The Von
Mises distribution already drives phases toward the forward cascade
direction. Blow-up cannot be prevented by phase alignment being
``unlikely''---the cascade naturally favors it. The prevention must
come from \emph{directional coherence degradation} (depolarization),
not from phase improbability.
\end{remark}

\subsection{Assembly time convergence}

\begin{proposition}[Finite assembly time]\label{prop:assembly}
In a shell model with ratio $\lambda > 1$, the total time for the
cascade to relay information from the Fano shell to all higher scales
is
\[
  T_{\mathrm{assemble}} = \frac{\tau_0}{1 - \lambda^{-2/3}},
\]
which is finite. For $\lambda = 2$:
$T_{\mathrm{assemble}} \approx 2.70\,\tau_0$.
\end{proposition}

\begin{proof}
The eddy turnover time at shell $n$ is $\tau_n = \tau_0\lambda^{-2n/3}$
(Kolmogorov). The total is a geometric series:
$\sum_{n=0}^\infty \tau_0 \lambda^{-2n/3}
= \tau_0/(1 - \lambda^{-2/3})$.
\end{proof}

\begin{remark}[Timing does not prevent blow-up]
The convergence of $T_{\mathrm{assemble}}$ means the cascade
reaches all scales in $O(\tau_0)$ time. Any argument against blow-up
must therefore be about \emph{amplitude} or \emph{coherence}, not
timing.
\end{remark}

\subsection{The cross-helical necessity chain}

The following chain of results localizes the blow-up question to
a specific sector of the nonlinearity.

\begin{theorem}[Biferale--Titi, 2013]\label{thm:BT}
Consider the helical-decimated NS equations on $\mathbb{T}^3$: the
dynamics restricted to triads of a single helicity sign
($s_k = s_p = s_q$). For initial data $v_0 \in \dot{H}^{1/2}$ and
forcing $f \in L^2([0,T]; H^{-1/2})$, the solution satisfies
\[
  v^+ \in L^\infty([0,T]; \dot{H}^{1/2}) \cap L^2([0,T]; \dot{H}^{3/2})
\]
for all $T > 0$. In particular, no blow-up occurs.
\end{theorem}

The mechanism: with same-helicity interactions only, helicity
$\mathcal{H} = \sum_k k |u_k^+|^2$ is a sign-definite inviscid
invariant equivalent to $\|v\|_{H^{1/2}}^2$. This gives the a~priori
bound that the full NS lacks~\cite{BiferaleTiti2013}.

\begin{proposition}[Cross-helical dominance of forward cascade]
\label{prop:cross-dominant}
Removing all same-helicity (Class~I) triads from the full NS
nonlinearity has no effect on the energy cascade or intermittency.
The forward cascade is entirely driven by cross-helical
interactions~\cite{SahooBiferale2018}.
\end{proposition}

\begin{corollary}[Blow-up requires cross-helical interactions]
\label{cor:blowup-cross}
Any finite-time singularity of the 3D NS equations must involve
cross-helical triadic interactions. The same-helical sector is globally
regular (Theorem~\ref{thm:BT}) and does not produce a forward cascade
(Proposition~\ref{prop:cross-dominant}).
\end{corollary}

\begin{proposition}[Leray suppression of the dangerous sector]
\label{prop:chain-suppression}
Cross-helical interactions are suppressed by the Leray projector.
The fraction of cross-helical Lamb vector surviving projection is
$\alpha(\theta) = (1-\cos\theta)/(3-\cos\theta)$
(Lemma~\ref{lem:leray}), with isotropic average
$\langle \alpha \rangle = 1 - \ln 2 \approx 0.307$. That is,
$\ln 2 \approx 69.3\%$ of the sole sector capable of driving blow-up
is killed by incompressibility.
\end{proposition}

\begin{remark}[The gap: suppression $\neq$ elimination]
The Biferale--Titi argument succeeds because removing cross-helical
triads \emph{entirely} makes helicity sign-definite, yielding
$\|v\|_{H^{1/2}}$ bounds. Our Leray suppression \emph{weakens}
cross-helical coupling by $69.3\%$ but does not eliminate it. The
surviving $30.7\%$ still permits sign-indefinite helicity contributions.
Whether this weakened coupling can drive blow-up depends on phase
coherence --- which is the subject of the directional bottleneck below.
\end{remark}

\subsection{The directional bottleneck conjecture}

We now combine all the above into a structural argument.

\begin{definition}[Directional information capacity]
Let $\mathcal{C}_{\mathrm{dir}}(K)$ denote the number of independent
directions in mode space at scale $K$ that carry coherent phase
information from the Fano shell through the cascade. Formally, this is
the effective rank of the mutual information
$I(X_{\mathrm{Fano}}; X_K)$ projected onto the mode eigenspaces at $K$.
\end{definition}

\begin{conjecture}[Directional Bottleneck]\label{conj:bottleneck}
Let $N(K) = \log_\lambda(K/k_F)$ denote the number of cascade steps
from the Fano shell to wavenumber $K$. Then:
\begin{enumerate}
  \item The directional capacity satisfies
    $\mathcal{C}_{\mathrm{dir}}(K) \leq 6 \cdot \rho^{N(K)}$
    where $\rho < 1$ is the depolarization factor per cascade step
    (Theorem~\ref{thm:depol}).
  \item Blow-up at scale $K$ requires coherent organization of
    $\sim K^2$ modes (shell surface), necessitating directional
    capacity $\mathcal{C}_{\mathrm{needed}}(K) \geq c\,K^\alpha$
    for some $c, \alpha > 0$.
  \item Since $6\rho^{N(K)}$ decays exponentially while
    $cK^\alpha$ grows polynomially,
    $\mathcal{C}_{\mathrm{dir}} < \mathcal{C}_{\mathrm{needed}}$
    for all $K$ sufficiently large. Phase-coherent blow-up is
    impossible.
\end{enumerate}
\end{conjecture}

The logic: the Fano shell provides 6 channels of information
(Proposition~\ref{prop:fano-lap}), all isotropic (maximal degeneracy
of $K_7$). Each cascade step depolarizes
(Theorem~\ref{thm:depol}), reducing the directional content
exponentially. After $O(1)$ steps, only scalar information (total
energy = $S_0$) survives. But blow-up requires \emph{specific}
phase alignment at small scales, which is directional information.
Energy arrives; coherence does not.

\subsection{Gaps and status}

Three mathematical gaps prevent Conjecture~\ref{conj:bottleneck} from
being a theorem:

\begin{enumerate}
  \item \textbf{Rigorous definition of $\mathcal{C}_{\mathrm{dir}}$.}
    The ``effective rank of projected mutual information'' needs a
    precise functional-analytic definition. The concept is new; no
    prior work decomposes cascade information flow by direction.

  \item \textbf{Can local dynamics self-organize?} The argument assumes
    all phase organization at scale $K$ must be relayed from the Fano
    shell. If local triadic interactions at scale $K$ can independently
    create coherent phase alignment, the bottleneck is bypassed.
    Three mechanisms oppose local self-organization
    (Waleffe null + Leray $\sin^2\phi/4$ + Buaria anti-twist), but
    their combined sufficiency is not proved. The per-mode balance of
    nonlinear self-organization $(\tfrac{1}{4}K|u_K|)$ against viscous
    damping $(\nu K^2)$ gives a crossover at $|u_K| = 4\nu K$: viscosity
    has the \emph{scale advantage} ($K^2$ vs.\ $K$), winning at all
    $K \gtrsim 0.35\,k_d$ (exact: $\sqrt{2}/4 \approx 0.354$) under
    Kolmogorov scaling. The $\sin^2\phi/4$ factor further lowers the
    effective crossover, though the precise value depends on the
    averaging convention for $\phi$ over triadic geometries.

  \item \textbf{Shell model $\to$ full 3D NS.} Tanogami--Araki results
    (\textsc{P1}--\textsc{P4}) are proved for shell models. The
    identity $\dot{\mathcal{I}}_K = \dot{\sigma}_K$ holds for any
    deterministic system (Liouville structure), but the specific bound
    $\dot{\mathcal{I}}_K \leq C \cdot K \cdot |u_K|$ uses shell-model
    coupling. Extension to full 3D NS requires bounding the triadic
    coupling tensor $B(k,p,q)$, where $\sin^2\phi/4$ from
    Lemma~\ref{lem:leray} enters as the structural cap.
\end{enumerate}

%% ========================================================
\section{The shape as entropy producer}
%% ========================================================

The preceding sections establish a chain of structural results:
the Fano plane constrains the cascade topology (Section~5),
depolarization destroys directional coherence (Section~5.4),
and the bottleneck conjecture bounds directional capacity
(Section~5.6). We now interpret these results physically
and propose quantitative tests.

\subsection{Three perpendicular pipes}

The Fano plane $\mathrm{PG}(2,2)$ is the projective plane over
$\mathbb{F}_2$. Its 7~points are the non-zero vectors of
$\mathbb{F}_2^3$, which on $\mathbb{T}^3$ correspond to integer
wavevectors:
\[
  \underbrace{(1,0,0),\; (0,1,0),\; (0,0,1)}_{\text{3 orthogonal axes}}
  ,\quad
  (1,1,0),\; (1,0,1),\; (0,1,1),\quad (1,1,1).
\]
The three coordinate directions are mutually perpendicular---not
by convention, but by the structure of $\mathbb{F}_2^3$ embedded
in Euclidean $\mathbb{R}^3$.

\begin{proposition}[Perpendicular pipes]\label{prop:perp-pipes}
Any flow sustaining all 7 triadic interactions of the Fano plane
must have significant energy in at least 3 mutually perpendicular
directions at the Fano shell $|k| \leq 3$.
\end{proposition}

\begin{proof}
The 7 lines of $\mathrm{PG}(2,2)$ are:
$\{e_1, e_2, e_1{+}e_2\}$,
$\{e_1, e_3, e_1{+}e_3\}$,
$\{e_2, e_3, e_2{+}e_3\}$,
$\{e_1, e_2{+}e_3, e_1{+}e_2{+}e_3\}$,
$\{e_2, e_1{+}e_3, e_1{+}e_2{+}e_3\}$,
$\{e_3, e_1{+}e_2, e_1{+}e_2{+}e_3\}$,
$\{e_1{+}e_2, e_1{+}e_3, e_2{+}e_3\}$.
Each basis vector $e_i$ appears in exactly 3 lines.
If any $\hat{u}_{e_i} = 0$, those 3 lines are inactive
(their triadic coupling vanishes).
Since each line involves exactly 3 modes summing to zero
in $\mathbb{F}_2^3$, inactivity of any basis direction
kills nearly half the cascade channels.
Thus all three basis directions---which are mutually
orthogonal---must carry energy for the full Fano
cascade to operate.
\end{proof}

\begin{remark}[Connection to local regularity]
Two perpendicular vortex tubes suffice to break cone
confinement: if $\omega$ has significant components along
two directions with angle $\theta > \pi/4$, no single
cone of half-angle $\pi/4$ contains both.
The Fano structure provides three such directions, exceeding
the local regularity threshold. This is a structural
obstruction to single-tube blow-up scenarios: the cascade
cannot operate through a single coherent tube.
\end{remark}

\subsection{The regularity dichotomy}

At any wavenumber~$K$ in the cascade, define the
\emph{directional diversity}
\[
  D(K) = \text{number of independent directions at scale } K
  \text{ carrying } \geq \delta \text{ fraction of the shell energy},
\]
for a fixed threshold $\delta > 0$.

\begin{proposition}[Dichotomy]\label{prop:dichotomy}
For any $K$ in the active cascade range, exactly one
of the following holds:
\begin{enumerate}
  \item[\textup{(A)}] \textbf{Multi-directional:}
    $D(K) \geq 2$. Then the dihedral angle $\beta > 0$
    between consecutive triadic planes, and enstrophy
    contributions accumulate as $\sqrt{N}$ (random walk),
    which is integrable.
  \item[\textup{(B)}] \textbf{Uni-directional:}
    $D(K) = 1$. Then the energy flux from the Fano shell
    to scale $K$ is bounded by
    $\dot{\mathcal{I}}_K \leq (7C/4) \cdot k_F \sqrt{E} / D(k_F)$,
    since only one of $D(k_F) \geq 3$ perpendicular channels
    is accessible. The reduced flux bounds the enstrophy
    production at scale $K$.
\end{enumerate}
In both cases, enstrophy growth at scale $K$ is controlled.
\end{proposition}

The physical content: blow-up requires both
\emph{coherent alignment} ($D = 1$, linear accumulation)
and \emph{sustained energy flux} (to feed the growing enstrophy).
The Fano structure makes these contradictory:
coherent alignment restricts the flow to one of three
perpendicular channels, cutting the available flux by
at least $2/3$.

\subsection{Fluid to heat: the cascade as entropy production}

The energy cascade converts organized kinetic energy
(coherent, directional, phase-structured) into
disorganized thermal energy (incoherent, isotropic, random).
The Fano plane is the throat of this transformation:

\begin{enumerate}
  \item \textbf{Input} (large scales, $|k| \leq 3$):
    Energy enters with spatial structure---vortices, jets,
    coherent motions. Phase relationships between modes
    encode this structure.

  \item \textbf{Processing} (Fano bottleneck):
    The $K_7$ Laplacian has maximally degenerate spectrum
    (Proposition~\ref{prop:fano-lap}): all 6 non-trivial
    eigenvalues equal $7/4$. No preferred diffusion direction.
    Information passes through 6 isotropic channels at equal rate.
    The 3 perpendicular pipes ensure no direction is favored.

  \item \textbf{Depolarization} (cascade relay):
    Each cascade step randomizes directional content
    (Theorem~\ref{thm:depol}): cross-helical dichroism~$= 0$,
    same-helical anti-polarizing. After $O(1)$ steps,
    only scalar information (total energy) survives.

  \item \textbf{Output} (dissipation scale, $k \sim k_d$):
    Energy arrives isotropic and phase-random.
    Viscosity converts it to heat.
    The transformation is complete.
\end{enumerate}

The regularity question, in this language, becomes:
\emph{is the entropy production rate always finite?}
The Fano throughput bound
(Proposition~\ref{prop:fano-throughput}) caps the processing
rate at $\dot{\mathcal{I}}_{\mathrm{Fano}}
\leq (7C/4)\, k_F \sqrt{E}$.
If this bound holds---if the throat cannot choke---then
energy cannot accumulate at any scale, and singularities
are impossible.

The three perpendicular pipes, the maximal isotropy,
and the Leray suppression all serve the same function:
they keep the throat open. Three independent channels
prevent blockage. Isotropy prevents preferential routing.
Leray suppression ($\alpha \approx 0.307$) reduces the
coherent fraction that could organize a blockage.

\subsection{Testable predictions}\label{sec:tests}

The entropy-production interpretation yields five
quantitative predictions, all testable in DNS or
shell-model simulations.

\medskip\noindent
\textbf{Test 1: Directional entropy at the Fano shell.}
Define the directional entropy at scale $K$:
\[
  S_{\mathrm{dir}}(K) = -\sum_{i=1}^{M}
  p_i(K) \ln p_i(K),
  \qquad
  p_i(K) = \frac{E_i(K)}{\sum_j E_j(K)},
\]
where $E_i(K)$ is the energy in direction $i$ at scale $K$,
summed over all modes with $|k| = K$ projecting onto
direction~$i$.
\emph{Prediction:} At the Fano shell ($K \leq 3$),
$S_{\mathrm{dir}} \geq \ln 3 \approx 1.10$
(energy distributed across 3 perpendicular pipes).
$S_{\mathrm{dir}}$ should not decrease with increasing~$K$
if depolarization holds.

\medskip\noindent
\textbf{Test 2: Phase correlation decay.}
Measure the phase correlation between Fano-shell modes
and modes at scale $K$:
\[
  C_{\mathrm{phase}}(K) =
  \left|\left\langle
    e^{i\theta_{k_F}} \, e^{-i\theta_K}
  \right\rangle\right|,
\]
averaged over mode pairs connected by the cascade.
\emph{Prediction:}
$C_{\mathrm{phase}}(K) \sim \rho^{N(K)}$
with $\rho < 1$ (exponential decorrelation).
If $C_{\mathrm{phase}}$ plateaus or increases at some scale,
local self-organization may be occurring
(violating the bottleneck assumption).

\medskip\noindent
\textbf{Test 3: Directional flux splitting.}
At the Fano shell, decompose the energy flux
$\Pi(k_F)$ into contributions from each of the
3 perpendicular directions.
\emph{Prediction:} Each direction carries
$(33 \pm 10)\%$ of the total flux (approximate
equipartition, consistent with maximal isotropy of $K_7$).
Deviations exceeding $50\%$ in a single direction
would indicate a preferred cascade channel,
weakening the scrambling argument.

\medskip\noindent
\textbf{Test 4: Suppression of uni-directional cascades.}
In DNS with controlled initial conditions, launch energy
along a single axis (e.g., Taylor--Green vortex along~$z$).
Monitor $D(K)$ as the cascade develops.
\emph{Prediction:} $D(K)$ increases from~1 to~$\geq 3$
within $O(1)$ large-eddy turnover times.
The Fano structure forces redistribution across
all 3 perpendicular pipes. If the flow remains
uni-directional for $t \gg \tau_0$, the Fano
cascade is not operating, and energy should
accumulate at large scales.

\medskip\noindent
\textbf{Test 5: Enstrophy accumulation statistics.}
Measure the scaling of $\omega \cdot S \cdot \omega$
contributions across $N$ consecutive triadic shells.
\emph{Prediction:} In the multi-directional regime
($D \geq 2$), the sum grows as $\sqrt{N}$.
In artificially constrained uni-directional flow,
the sum grows as $N$ (linear).
The transition between regimes should occur at
$D \approx 2$, consistent with the dichotomy
(Proposition~\ref{prop:dichotomy}).

\bigskip
Tests 1--3 require standard DNS post-processing
(energy spectra, phase statistics, flux decomposition).
Test~4 requires only a Taylor--Green initial condition,
which is standard. Test~5 requires tracking enstrophy
contributions shell-by-shell, which is more demanding
but feasible at moderate Reynolds numbers
($\mathrm{Re} \leq 1600$).

%% ========================================================
\section{What remains}
%% ========================================================

\subsection{The verified structure}

The following results are independent of the curvature error and
stand on their own:

\begin{enumerate}
  \item The inner product $(\cos\phi + \sigma\tau)/2$ generates the
    helical structure of triadic interactions.
  \item The Waleffe null kills same-helical local triads
    (kinematic, exact).
  \item Leray suppression kills 69.3\% of cross-helical Lamb vectors
    (kinematic, exact).
  \item Equipartition: both channels have Leray-weighted coupling
    $1/6$ (geometric, exact).
  \item Numerically: cross-helical dominance in cancellation is
    confirmed on NS-evolved fields.
  \item Numerically: cancellation increases with Re (two data points).
  \item Arnold curvature is uniformly negative and
    helicity-independent (Araki 2016, independently verified).
  \item Cross-helical necessity chain: blow-up requires cross-helical
    triads (Biferale--Titi + Sahoo--Biferale), which are suppressed
    69.3\% by Leray.
  \item Universal Depolarization: cross-helical dichroism $= 0$,
    same-helical anti-polarizing (symbolic + numerical).
  \item Fano Laplacian: algebraic connectivity $\lambda_2 = 7/4$
    with maximal degeneracy (algebraic, exact).
  \item Von Mises phase cost: $-\ln p(\pi/2) \to
    \tfrac{1}{2}\ln(2\pi/\kappa)$ (negative for $\kappa \gtrsim 6.55$:
    cascade naturally favors blow-up phases).
  \item Assembly time finite: $T_{\mathrm{assemble}} \approx
    2.70\,\tau_0$ for $\lambda = 2$ (geometric series).
  \item Perpendicular pipes: $\mathbb{F}_2^3$ embedding forces
    3 mutually orthogonal directions at the Fano shell (algebraic, exact).
  \item Regularity dichotomy: multi-directional $\Rightarrow$ $\sqrt{N}$
    accumulation; uni-directional $\Rightarrow$ energy starvation
    (structural, from Fano + throughput bound).
\end{enumerate}

\subsection{Open questions}

\begin{enumerate}
  \item \textbf{Directional information capacity.}
    Formalize $\mathcal{C}_{\mathrm{dir}}(K)$ and prove the
    exponential decay through depolarization.
  \item \textbf{Local self-organization.} Can local triadic
    interactions at scale $K$ create blow-up-relevant phase coherence
    without relay from the Fano shell? Anti-twist opposes it; proof
    is lacking.
  \item \textbf{Shell model $\to$ 3D NS.} Extend the Tanogami--Araki
    information flow bounds from shell models to full 3D
    Navier--Stokes.
  \item \textbf{Does cancellation scale as $K^{-\alpha}$ with
    $\alpha > 1$?} Two Re data points show improving cancellation.
    Whether $|dZ/dt|/Z^{3/2}$ stays bounded as $\mathrm{Re} \to \infty$
    is the Millennium Problem.
  \item \textbf{Quadrature bridge to Constantin--Fefferman.}
    The direction change $\delta\xi(k) \sim \tfrac{1}{4} R_K(k)
    \cdot k^{1/6}$ from shell~$k$ (under Kolmogorov scaling) gives
    the arithmetic bridge
    $|\xi(x)-\xi(y)| \leq \tfrac{1}{4}\sum_k R_K(k)\, k^{1/6}\,
    \min(1, k\rho)$,
    requiring $q > 7/6$ for convergence.
    If inter-shell contributions are \emph{independent}
    (justified by: (i)~$\sin^2\phi/4$ selects orthogonal triads
    whose rotation axes are geometrically unrelated across scales;
    (ii)~Berry frustration, Chern number~2, prevents global phase
    alignment), the sum becomes quadrature:
    $|\xi(x)-\xi(y)|^2 \leq \tfrac{1}{16}\sum_k R_K^2(k)\, k^{1/3}\,
    \min(1, k^2\rho^2)$,
    lowering the convergence threshold to $q > 2/3$
    while the H\"older-$1/2$ threshold remains $q \geq 7/6$.
    DNS at $\mathrm{Re} = 1600$ gives $q_{\mathrm{local}} \approx 1.17$,
    yielding exactly H\"older-$1/2$ under quadrature.
    Whether sub-H\"older-$1/2$ direction regularity
    ($q \in (2/3, 7/6)$, corresponding to
    H\"older-$(q - 2/3) < 1/2$) suffices for global regularity
    is an open problem in PDE theory~\cite{BdV2019}.
\end{enumerate}

\begin{remark}[Kolmogorov duality]
The threshold $q = 7/6$ is not accidental. It equals
$1 + \tfrac{1}{6}$, where $k^{1/6}$ is the shell vorticity amplitude
under Kolmogorov scaling $E(k) \sim k^{-5/3}$.
Thus $q = 7/6$ is the phase decoherence exponent that
\emph{exactly compensates} vorticity growth across scales:
the bridge sum converges if and only if decoherence
outpaces the $k^{1/6}$ amplitude factor.

This is dual to the Kolmogorov cascade itself.
If $q < 7/6$, vorticity direction becomes irregular,
stretching is uncontrolled, energy accumulates at small scales
(spectrum steeper than $k^{-5/3}$), dissipation increases,
and the cascade restores itself.
If $q > 7/6$, direction is too smooth, stretching too weak,
energy does not cascade efficiently (shallower spectrum),
and turbulence intensifies until $q$ decreases.
The Kolmogorov state $E(k) \sim k^{-5/3}$ and the
regularity threshold $q = 7/6$ are the same fixed point,
expressed in dual languages.

The feedback is quantitative: for a general spectrum
$E(k) \sim k^{-p}$, the convergence threshold is
$q^*(p) = 2 - p/2$, with $dq^*/dp = -1/2$.
As the spectrum steepens ($p$ increases toward a concentration
event), the required decoherence \emph{drops}.
At $p = 2$: $q^* = 1$.
At $p = 3$ (enstrophy cascade): $q^* = 1/2$.
The system fights its own blow-up---energy concentration
makes regularity \emph{easier} to maintain.

The Millennium Problem, in this formulation, asks:
during a concentration event, does the actual phase
decoherence $q(p)$ decrease faster or slower than the
threshold $q^*(p) = 2 - p/2$?
If $dq/dp > -1/2$, the regularity margin $q - q^*$
\emph{grows} under concentration, and blow-up is
self-defeating. If $dq/dp < -1/2$, coherent structures
created by concentration destroy decoherence faster than
the threshold drops, and blow-up can escape the feedback.
\end{remark}

\bigskip
\begin{center}\rule{2in}{0.4pt}\end{center}

\medskip
\noindent\textit{Seven points of a cube with one corner missing.
The missing corner is the ground. The seven that remain form the Fano
plane. Every pair meets on exactly one line. Every line is a triad.
The floor is a sphere. The crack is $3/28$. More thread than tear.}

\begin{thebibliography}{99}
\bibitem{Arnold1966} V.\,I.~Arnold, Sur la g\'eom\'etrie
  diff\'erentielle des groupes de Lie de dimension infinie et ses
  applications \`a l'hydrodynamique des fluides parfaits,
  \textit{Ann.\ Inst.\ Fourier} \textbf{16} (1966), 319--361.

\bibitem{Waleffe1992} F.~Waleffe, The nature of triad interactions
  in homogeneous turbulence, \textit{Phys.\ Fluids A} \textbf{4}
  (1992), 350--363.

\bibitem{Araki2016} R.~Araki, Sectional curvature of the
  volume-preserving diffeomorphism group of $\mathbb{T}^3$ and
  complex helical wave modes, preprint (2016), arXiv:1608.05154.

\bibitem{ConstantinFefferman1993} P.~Constantin, C.~Fefferman,
  Direction of vorticity and the problem of global regularity for the
  Navier--Stokes equations, \textit{Indiana Univ.\ Math.\ J.}
  \textbf{42} (1993), 775--789.

\bibitem{Miller2024} E.~Miller, A model equation for enstrophy
  growth and its global regularity, preprint (2024), arXiv:2407.02691.

\bibitem{TanogamiAraki2024} R.~Tanogami, R.~Araki,
  Information-thermodynamic bound on information flow in turbulent
  cascade, \textit{Phys.\ Rev.\ Research} \textbf{6} (2024), 013090;
  arXiv:2206.11163.

\bibitem{TanogamiAraki2025} R.~Tanogami, R.~Araki, Scale-to-scale
  information flow amplifies turbulent fluctuations,
  \textit{Phys.\ Rev.\ Research} \textbf{7} (2025), 023078;
  arXiv:2408.03635.

\bibitem{Tanogami2025} R.~Tanogami, Scale locality of information
  flow in turbulence, \textit{J.\ Stat.\ Mech.} (2025), 093405;
  arXiv:2407.20572.

\bibitem{BenavidesBustamante2026} S.\,J.~Benavides,
  M.\,D.~Bustamante, Phase dynamics and energy flux in hydrodynamic
  shell models, preprint (2025, revised 2026), arXiv:2507.03397.

\bibitem{Buaria2024} D.~Buaria, J.\,M.~Lawson, M.~Wilczek,
  Anti-twist and self-regularization in three-dimensional
  vortex dynamics, \textit{Sci.\ Adv.} \textbf{10} (2024), eado1969.

\bibitem{PecoraCarroll1998} L.\,M.~Pecora, T.\,L.~Carroll,
  Master stability functions for synchronized coupled systems,
  \textit{Phys.\ Rev.\ Lett.} \textbf{80} (1998), 2109--2112.

\bibitem{BiferaleTiti2013} L.~Biferale, E.\,S.~Titi,
  On the global regularity of a helical-decimated version of the 3D
  Navier--Stokes equations, \textit{J.\ Stat.\ Phys.}
  \textbf{151} (2013), 1089--1098; arXiv:1303.1215.

\bibitem{SahooBiferale2018} G.~Sahoo, L.~Biferale,
  Disentangling the triadic interactions in Navier--Stokes equations,
  \textit{Fluid Dyn.\ Res.} \textbf{50} (2018), 011404;
  arXiv:1709.03713.

\bibitem{BdVBerselli2002} H.~Beir\~ao da Veiga, L.\,C.~Berselli,
  On the regularizing effect of the vorticity direction in
  incompressible viscous flows, \textit{Differential Integral
  Equations} \textbf{15} (2002), 345--356.

\bibitem{BdV2019} H.~Beir\~ao da Veiga,
  Navier--Stokes equations: some questions related to the direction
  of the vorticity, \textit{Discrete Contin.\ Dyn.\ Syst.\ Ser.\ S}
  \textbf{12} (2019), 203--213.

\bibitem{GrujicRuzmaikina2004} Z.~Gruji\'c, A.~Ruzmaikina,
  Interpolation between algebraic and geometric conditions for
  smoothness of the vorticity in the 3D NSE, \textit{Indiana Univ.\
  Math.\ J.} \textbf{53} (2004), 1073--1080.
\end{thebibliography}

\end{document}
